\begin{frame}{보고서 구조 (M4) — 발표 슬라이드 = 보고서의 절 제목}
  \small
  \begin{enumerate}
    \item \textbf{문제 정의} (0.5p): MDP 5요소(Ch3.1), 상태·행동 차원과
      범위, 전이·보상(이론적 리턴 범위 포함), 할인율, \textbf{주 지표}
      (학습 후 결정적 평가의 평균 리턴), \textbf{베이스라인}.
    \item \textbf{알고리즘 선택과 이유} (1p): \textbf{왜 이 환경을 이
      알고리즘에 골랐는가}(§페어링의 ``무대'' 논리), 하이퍼파라미터
      ($\gamma, \lambda, \epsilon_{\text{clip}}$, 보상 스케일)와 선택
      근거(Ch10$\sim$13 개념).
    \item \textbf{실험 프로토콜} (1p): 시드 목록, 학습 예산,
      \textbf{평가 프로토콜(결정적 $a=\mu$, 에피소드 수, 모든 시드에 동일)
      명시}, 하이퍼파라미터 선택 절차. \textbf{이 절이 없으면 16.2
      체크리스트에서 바로 감점.}
    \item \textbf{결과} (1p): 시드 비교 표, 학습 곡선 1장(여러 시드를
      얇은 선으로), \textbf{``같은 정책 여러 측정'' 패널}(학습
      도중/결정적/베이스라인).
    \item \textbf{수식적 정당화} (1p): Ch9$\sim$15 개념 $\geq 1$개로
      ``왜 이 결과가 나왔는지'' — ``좋은 예'' 문체로.
    \item \textbf{잘 안 된 부분과 원인 진단} (0.5p): \textbf{필수} —
      안 된 부분이 없었으면 ``무엇을 시도했다가 안 됐는가''를 쓴다.
      ``모두 잘 됐다''는 보고서는 이 절이 비어 있다는 뜻.
  \end{enumerate}
\end{frame}
